فرض کنید تصویر اولیه $f(x,y)$ حقیقی است. در گام اول داریم
$$g(x,y)=f(x,y)(-1)^{x+y}.$$
اگر تبدیل فوریه دوبعدی آن را با $G(u,v)$ نشان دهیم، پس
$$G(u,v)=\mathcal{F}\{g(x,y)\}.$$
گام سوم گرفتن مزدوج مختلط است:
$$\tilde{G}(u,v)=G^{*}(u,v).$$
حال تبدیل معکوس می‌دهد
$$\tilde{g}(x,y)=\mathcal{F}^{-1}\{\tilde{G}(u,v)\}.$$

برای تبدیل گسسته دوبعدی داریم
$$\mathcal{F}^{-1}\{G^{*}(u,v)\}=\tilde{g}(x,y)=g^{*}(-x,-y).$$
چون تصویر حقیقی است، $g^{*}(-x,-y)=g(-x,-y)$.

در گام آخر، شیفت مکانی را برمی‌گردانیم:
$$\tilde{f}(x,y)=\tilde{g}(x,y)(-1)^{x+y}=g(-x,-y)(-1)^{x+y}.$$
اما
$$g(-x,-y)=f(-x,-y)(-1)^{-x-y}=f(-x,-y)(-1)^{x+y}.$$
پس نتیجه:
$$\tilde{f}(x,y)=f(-x,-y).$$

یعنی تصویر خروجی چرخش $180^\circ$ (تقارن مرکزی) تصویر ورودی است؛ به همین دلیل شکل (b) همان شکل (a) با دوران $180^\circ$ دیده می‌شود.